\documentclass[letterpaper,11pt]{article}

\usepackage[empty]{fullpage}
\usepackage[T1]{fontenc}
\usepackage{mathptmx}
\usepackage[colorlinks=true, linkcolor=NavyBlue, urlcolor=NavyBlue]{hyperref}
\usepackage{parskip}

\addtolength{\oddsidemargin}{-0.6in}
\addtolength{\evensidemargin}{-0.6in}
\addtolength{\textwidth}{1.2in}
\addtolength{\topmargin}{-0.7in}
\addtolength{\textheight}{1.6in}

\setlength{\parskip}{5pt}

\begin{document}

\noindent\textbf{\large Charles Benello} \hfill \today \\
\noindent +1(312) 989-9950 $|$ cmbenello@uchicago.edu $|$ Chicago, IL

\vspace{6pt}

\noindent\textbf{Exponential Technologies -- Quantitative Research Summer Intern}

\vspace{2pt}

I'm a UChicago Master's student in Financial Mathematics applying for the Quantitative Research Summer Intern role. What pulled me toward Exponential is Tesseract Agentic Analytics: the combination of institutional flow data, macro forecasting, and an agentic interface on top is unusually close to the kind of work I've been gravitating toward over the last year.

The role splits across six tracks. I'd be happy on any of them, but I want to be specific about where my background is strongest and where I'd be ramping, so the staffing decision can be driven by fit rather than my preference.

\textbf{Strongest fit: AI-Driven Backtesting Agent.} At Bank of America (UChicago Project Lab) last fall I built an AI-driven platform to rebalance HQLA portfolios under capital, liquidity, and net interest income constraints. The core was a multi-agent LLM debate system: competing agents generate macro stress scenarios, debate the implications, and produce probability-weighted reallocation recommendations. The architecture is the same shape as what your agent project describes -- users defining and stress-testing strategy components, with an agent orchestrating the experiment loop. I shipped a Next.js scenario explorer, a Flask yield curve dashboard, a FastAPI service exposing the optimizer, and a daily ingestion pipeline pulling UST yield curves and macro data into structured snapshots.

Backtesting an agent is fundamentally an evaluation problem, which is the half I work on at my current SWE internship at Alter Domus. I built an LLM-as-a-Judge framework with automated benchmarks for faithfulness, numerical accuracy, and completeness. The same patterns -- scoring against ground truth, isolating modes of failure, regression-testing prompts and configs -- carry over to evaluating an agent that defines and tests strategy components.

\textbf{Time-series modeling at scale: PayPal foundation model.} I'm currently a quantitative researcher at PayPal through Project Lab, building a self-supervised transaction foundation model: a TabBERT-style hierarchical 2-level transformer ($\sim$5M params) on 14.5M unlabeled transactions, evaluated on the IBM TabFormer benchmark (24M txns, 0.12\% fraud rate). I've run 87+ pretraining ablations on UChicago's Midway3 SLURM cluster comparing BERT-style vs.\ GPT-style architectures, with sweeps across tokenization, time encoding (score-relative vs.\ inter-event), masking rate, and sequence length. I implemented per-field masking from PRAGMA (Revolut 2026) and the SPRM module from PANTHER (NeurIPS 2025), with per-field masking alone yielding +1.9\% AUC and +16.8\% F1. Best result so far is AUC 0.974 / F1 0.822. Most of this is transformer-based time-series rather than ARIMA-style classical econometrics, but I've spent real time on temporal encodings and probability calibration on heavily imbalanced data, which is the regime where flow-signal evaluation also lives.

\textbf{White papers.} Three published papers in database systems and resource-adaptive query execution (VLDB 2026, CIDR 2025, ADMS 2026), one sole-authored. Comfortable writing at academic length and have a sense of how to scope a paper from problem framing through publication.

\textbf{Honest about limitations.} My ML background leans neural and self-supervised. I've used XGBoost and MLPs as baselines at PayPal, but I'm not coming in fluent on tree-based ML or classical ARIMA-style econometrics. For the macro-forecast and flow-analytics tracks specifically, I'd be ramping up on those tools rather than leading from day one. That's a fast ramp, but I'd rather be upfront about it.

The unpaid structure is fine on my end. What I'm looking for is a chance to ship real work with people who already think carefully about institutional flow data and signal development. I'm available the full summer and can hit the 6-hour core-team overlap comfortably from Chicago. Happy to share more detail on any of the above.

\vspace{4pt}

\noindent Charles Benello

\end{document}
